Teman Tuli atau difabel pendengaran di Indonesia umumnya menggunakan bahasa isyarat sebagai alat komunikasi utama dalam kehidupan sehari-hari.
Bahasa isyarat berperan penting sebagai sarana penyampaian makna, ekspresi, dan interaksi sosial dalam komunitas Tuli.
Keberadaan bahasa isyarat juga diakui secara internasional sebagai bentuk komunikasi nonverbal yang setara dengan bahasa lisan \cite{verushka2024signlanguages}.

Di Indonesia, terdapat dua sistem bahasa isyarat yang dikenal luas, yaitu Sistem Bahasa Isyarat Indonesia (SIBI) dan Bahasa Isyarat Indonesia (BISINDO).
SIBI disusun secara formal oleh pemerintah dengan mengadaptasi struktur bahasa isyarat asing dan tata bahasa Indonesia lisan.
Penggunaan SIBI lebih umum dijumpai dalam konteks formal dan institusional \cite{saraswati2022bisindo}.
BISINDO berkembang secara alami dalam komunitas Tuli dan digunakan secara luas dalam percakapan sehari-hari.
Struktur BISINDO bersifat lebih kontekstual dan tidak selalu mengikuti tata bahasa Indonesia secara langsung \cite{nilawaty2020sibi, bentara2024sibi-bisindo}.

Perbedaan karakteristik tersebut menjadikan BISINDO sebagai bahasa yang sangat bergantung pada gerakan tangan, orientasi tubuh, serta ekspresi non-manual.
Pola komunikasi visual tersebut membentuk ciri khas BISINDO sebagai bahasa yang kaya akan informasi gestural.
Pemahaman BISINDO masih menjadi tantangan bagi masyarakat umum.
Keterbatasan akses dan rendahnya pemahaman bahasa isyarat menyebabkan hambatan komunikasi antara Teman Tuli dan masyarakat luas \cite{lavenia2023tuliindonesia, arlinta2020keterbatasan}.

Perkembangan teknologi \textit{computer vision} dan kecerdasan buatan membuka peluang pemanfaatan teknologi sebagai sarana pendukung komunikasi bahasa isyarat.
Berbagai penelitian sebelumnya menunjukkan bahwa pendekatan berbasis pembelajaran mesin mampu mempelajari pola visual dari gerakan manusia.
Pendekatan tersebut memberikan potensi untuk mengembangkan sistem bantu penerjemahan bahasa isyarat berbasis visual.

Sebagian besar penelitian pengenalan bahasa isyarat berfokus pada pemodelan temporal menggunakan data berurutan.
Dominasi pendekatan temporal menimbulkan keterbatasan dalam pemahaman potensi pendekatan alternatif yang lebih ringan dan berbasis spasial.
Kajian mengenai kemampuan model spasial dua dimensi dalam menangkap pola gerakan bahasa isyarat masih relatif terbatas.
Kondisi tersebut membuka peluang penelitian untuk mengeksplorasi pemanfaatan model spasial berbasis \textit{Convolutional Neural Network} dua dimensi, khususnya arsitektur ResNet-2D, dengan representasi data berupa \textit{landmark} tubuh manusia.

Penelitian ini difokuskan pada pengenalan pola gerakan Bahasa Isyarat Indonesia berbasis \textit{landmark} tanpa mengabaikan konteks sosial dan etika penggunaan BISINDO.
Pendekatan yang diusulkan diharapkan dapat memberikan kontribusi akademik dalam pengembangan teknologi pendukung yang selaras dengan praktik dan nilai komunitas Tuli.
